\section*{Bab \ref{ch:b2:unicellular-diversity} --- Keanekaragaman Organisme Bersel Tunggal}

\begin{solution}{exo:b2:unicellular-diversity:1}
\omterm{def:b2:unicellular-diversity:unicellular}{Bersel tunggal}: satu sel menjalankan semua fungsi dan hidup sendiri;
\omterm{def:b2:unicellular-diversity:unicellular}{berkoloni}: sel yang serupa tetap melekat tetapi masing-masing dapat
hidup sendiri; \omterm{def:b2:unicellular-diversity:unicellular}{bersel banyak}: beberapa tipe sel, saling bergantung,
dan hanya garis nutfahnya yang berbiak. \emph{E.~coli} dan ragi
adalah \omterm{def:b2:unicellular-diversity:unicellular}{bersel tunggal} (tunasnya melepaskan diri lalu hidup sendiri);
\emph{Gonium} \omterm{def:b2:unicellular-diversity:unicellular}{berkoloni} (16 sel serupa, masing-masing sanggup
mendirikan koloni); \emph{Volvox} \omterm{def:b2:unicellular-diversity:unicellular}{bersel banyak} (sel soma yang fana,
beberapa sel nutfah).
\end{solution}

\begin{solution}{exo:b2:unicellular-diversity:2}
$S/V = 3/r$: \qty{6}{\micro m^{-1}} bagi kokusnya,
\qty{0.06}{\micro m^{-1}} bagi \omterm{def:b2:unicellular-diversity:protist}{protistanya}, selisih seratus kali lipat.
Kokusnya bernapas lebih cepat tiap satuan volume: setiap bagian
sitoplasmanya berjarak kurang dari satu mikrometer dari permukaan
tempat oksigen dan makanannya masuk, sedangkan bagian dalam \omterm{def:b2:unicellular-diversity:protist}{protistanya}
disuapi lewat permukaan seratus kali lebih kecil tiap satuan volume.
\end{solution}

\begin{solution}{exo:b2:unicellular-diversity:3}
Tanpa nukleus (sebuah \omterm{def:b2:unicellular-diversity:prokaryote}{nukleoid}) lawan selubung inti; tanpa mitokondria
maupun endomembran, dengan rantai respirasinya duduk pada membran
plasma, lawan organel; sebuah flagelum berupa heliks kaku berputar dari
flagelin yang diputar proton, lawan silia $9+2$ yang membengkok oleh
dinein dan ATP; juga ukurannya (\qty{1}{\micro m} lawan puluhan) dan
dinding peptidoglikannya. Sebuah bakteri tidak dapat berfagositosis:
dindingnya melarangnya menelan butiran, yang dikerjakan amuba dengan
pseudopodiumnya.
\end{solution}

\begin{solution}{exo:b2:unicellular-diversity:4}
Silia: gerak dan penyapuan makanan menuju mulutnya; alur mulut:
mencorongkan butiran ke kerongkongannya; vakuola makanan: mencerna
bakteri yang tertelan dengan enzim lisosom; \omterm{prop:b2:unicellular-diversity:osmoregulation}{vakuola kontraktil}:
mengeluarkan air yang masuk secara osmosis; makronukleus: nukleus yang
bekerja, ratusan salinan gennya ditranskripsi bagi kehidupan
sehari-hari; mikronukleus: nukleus diploid yang bersifat nutfah,
dipakai dalam meiosis dan konjugasi.
\end{solution}

\begin{solution}{exo:b2:unicellular-diversity:5}
$t = x^2/2D$: $(10^{-6})^2/10^{-9} = \qty{1}{ms}$ bagi bakterinya;
$(5\times 10^{-4})^2/10^{-9} = \qty{250}{s}$, kira-kira empat menit,
bagi amubanya. Amubanya tidak dapat mengandalkan difusi untuk
menyebarkan metabolitnya: ia membutuhkan aliran sitoplasma dan
pengangkutan bermotor sepanjang sitoskeletonnya, yang tidak dibutuhkan
bakterinya.
\end{solution}

\begin{solution}{exo:b2:unicellular-diversity:6}
$Re = \rho vL/\eta$: \emph{Paramecium} $10^3\times 10^{-3}\times
2\times 10^{-4}/10^{-3} = 0.2$; bakteri $10^3\times 2\times
10^{-5}\times 2\times 10^{-6}/10^{-3} = 4\times 10^{-5}$; perenang
$10^3\times 1\times 2/10^{-3} = 2\times 10^6$. Pada $Re \gg 1$
kelembaman membawa perenangnya terus sesudah kayuhan; pada
$Re \ll 1$ seretan kental menghentikan bakterinya dalam jarak satu
nanometer, sehingga ia bergerak hanya selama ia berdenyut.
\end{solution}

\begin{solution}{exo:b2:unicellular-diversity:7}
$v_s = 2r^2\Delta\rho g/9\eta = 2\times 10^{-10}\times 100\times
9.81/(9\times 10^{-3}) = \qty{2.2e-5}{m/s}$; \qty{50}{m} memakan
$50/2.2\times 10^{-5} = \qty{2.3e6}{s}$, 27 hari. Memerohkan
jari-jarinya membagi $v_s$ dengan 4: 108 hari. Menurunkan $\Delta\rho$
menjadi \qty{20}{kg/m^3} membaginya dengan 5: 135 hari.
\end{solution}

\begin{solution}{exo:b2:unicellular-diversity:8}
$V = \tfrac{4}{3}\pi\times 100\times 25\times 25 =
\qty{2.6e5}{\micro m^3}$. Dalam \qty{1800}{s}: $2.6\times 10^5/1800 =
\qty{145}{\micro m^3/s}$; karena $\qty{1}{\micro m^3} =
\qty{1e-15}{L}$, ini setara \qty{1.5e-13}{L/s}.
\end{solution}

\begin{solution}{exo:b2:unicellular-diversity:9}
Pelipatduaan sepanjang \qty{1}{mm} adalah kenaikan kira-kira
\qty{0.07}{\%} tiap mikrometer (karena $2^{1/1000} - 1 = 0.0007$),
jadi \qty{0.07}{\%} sepanjang selnya. Sebuah lari \qty{1}{s} menempuh
\qty{20}{\micro m}: \qty{1.4}{\%}. Untuk memilah selisih
\qty{0.07}{\%}, selnya harus mencacah sekitar $(1/0.0007)^2 = 2\times 10^6$ molekul
pada tiap ujungnya, jauh lebih banyak daripada yang
terikat pada beberapa ribu reseptornya dalam satu sekon; \qty{1.4}{\%}
menuntut sekitar \num{5000}, dan itu terjangkau. Perbandingan waktu
mengubah pengukuran ruang yang mustahil menjadi pengukuran yang
mungkin, dengan membiarkan larinya yang mengintegralkan.
\end{solution}

\begin{solution}{exo:b2:unicellular-diversity:10}
Permukaan $4\pi r^2$ dengan $r = \qty{250}{\micro m}$: \qty{7.9e5}{\micro
m^2}; volume sitoplasmanya $\approx$ permukaan $\times$ \qty{2}{\micro m}
$= \qty{1.6e6}{\micro m^3}$; rasionya \qty{0.5}{\micro m^{-1}}. Bagi
\emph{E.~coli}, sebuah silinder: $S/V = 2\pi r L/\pi r^2 L = 2/r =
\qty{4}{\micro m^{-1}}$. Si raksasa hanya delapan kali lebih buruk
daripada \emph{E.~coli}, bukan 500 kali, karena setiap titik
sitoplasmanya berjarak kurang dari \qty{2}{\micro m} dari permukaannya.
Vakuolanya menyimpan nitrat, yaitu penerima elektronnya, sehingga ia
dapat merespirasikan sulfida di dalam endapan tanpa oksigen di antara
pengadukan langka yang membawa nitrat
(\cref{ch:b2:microbial-metabolism}).
\end{solution}

\begin{solution}{exo:b2:unicellular-diversity:11}
Sebuah sel tidak dapat berenang dan membelah sekaligus. Koloni berisi
16 sel yang berhenti berenang selama berjam-jam pembelahannya
tenggelam, menurut \omterm{thm:b2:unicellular-diversity:stokes}{hukum Stokes}, sejauh jarak yang dapat diabaikan
(jari-jarinya beberapa puluh mikrometer, kelajuan pengendapannya
beberapa mikrometer tiap sekon); seluruh koloninya dapat membelah
tanpa kehilangan apa pun. Bola berisi \num{10000} sel berjari-jari
\qty{500}{\micro m} akan tenggelam seratus kali lebih cepat ---
beberapa milimeter tiap sekon, beberapa meter tiap jam --- ke luar dari
cahayanya: sepadan menjaga sebagian besar selnya tetap berenang
sementara beberapa membelah. Sel nutfahnya besar karena ia harus
mengusung sumber daya untuk membangun seluruh bola anak berisi ribuan
sel lewat pembelahan berturut-turut tanpa makan, sehingga ia dibekali
somanya.
\end{solution}

\begin{solution}{exo:b2:unicellular-diversity:12}
Sebuah klad adalah satu moyang beserta semua keturunannya; sebuah
tingkatan adalah taraf pengorganisasian yang dicapai beberapa garis
keturunan. \omterm{def:b2:unicellular-diversity:protist}{Protista} --- alga hijau (segaris dengan tumbuhan darat),
diatom (dengan alga cokelat), siliata (dengan dinoflagelata dan
\emph{Plasmodium}), amuba (dekat dengan hewan dan fungi) --- hanya
berbagi sel eukariot yang hidup sendiri, dan moyang bersama
terakhirnya adalah moyang seluruh eukariot, termasuk kita: sebuah
tingkatan. ``\omterm{def:b2:unicellular-diversity:prokaryote}{Prokariot}'' menyatukan bakteri dan arkea berdasarkan apa
yang tidak dimilikinya; arkea lebih dekat kepada eukariot daripada
kepada bakteri, sehingga istilah itu pun sebuah tingkatan.
Sianobakteri berasal dari satu moyang yang menemukan fotosintesis
beroksigen dan mencakup semua keturunannya (kecuali garis
kloroplasnya): sebuah klad.
\end{solution}

\begin{solution}{pb:b2:unicellular-diversity:1}
\textbf{1.} $1\times 1\times 0.1 = \qty{0.1}{mm^3} = \qty{0.1}{\micro L}$.
\textbf{2.} $24/\qty{0.1}{\micro L} = 240$ tiap \unit{\micro L} $=
\qty{2.4e5}{}$ sel tiap mililiter.
\textbf{3.} $V = \tfrac{4}{3}\pi\times 125 = \qty{524}{\micro m^3}$;
$S/V = 3/5 = \qty{0.6}{\micro m^{-1}}$.
\textbf{4.} $2.4\times 10^5\times 524 = \qty{1.26e8}{\micro m^3}$ tiap
mililiter, dan $\qty{1}{mL} = \qty{1e12}{\micro m^3}$: bagian sebesar
$1.3\times 10^{-4}$, yaitu \qty{0.013}{\%}.
\textbf{5.} $\qty{1000}{m^3} = \qty{1e9}{mL}$: $2.4\times 10^{14}$
sel.
\textbf{6.} Sianobakterinya adalah filamen sel yang jauh lebih kecil
(beberapa mikrometer), berwarna biru kehijauan dan tanpa kloroplas
maupun nukleus yang jelas, serta tidak berenang dengan flagelum;
alganya adalah satu sel bulat telur dengan dua flagelum, sebuah bintik
mata, dan satu kloroplas berbentuk cawan.
\textbf{7.} $384/24 = 16 = 2^4$: empat kali pelipatduaan dalam
\qty{48}{h}, $T = \qty{12}{h}$.
\textbf{8.} $k = \ln 2/T = 0.693/12 = \qty{0.058}{h^{-1}}$.
\textbf{9.} $t = \ln(10^8/2.4\times 10^5)/k = \ln(417)/0.058 =
6.03/0.058 = \qty{104}{h}$, kira-kira 4,3 hari.
\textbf{10.} Haranya (fosfat, nitrat) habis; sel-selnya saling
menaungi sehingga cahaya tiap selnya turun; pemangsanya (siliata,
rotifera, \emph{Daphnia}) berlipat ganda memakannya; pada
\qty{1e8}{} tiap mililiter sel-selnya akan mengisi kira-kira lima
persen volumenya dan menghabiskan $\mathrm{CO_2}$ terlarutnya.
\textbf{11.} Tumbuh hanya separuh waktunya: purata lajunya menjadi
separuh, \qty{208}{h}, yaitu 8,7 hari.
\textbf{12.} Cacahnya mengukur populasinya, bukan sel perorangan: jika
sebuah sel melepaskan 8 anak tiap \qty{36}{h}, populasinya dikalikan
8 $= 2^3$ dalam tiga \omterm{def:b2:unicellular-diversity:division}{waktu generasi} masing-masing \qty{12}{h}. \omterm{def:b2:unicellular-diversity:division}{Waktu
generasinya} adalah purata waktu pelipatduaan cacahnya, bagaimanapun
pembelahannya dikelompokkan.
\textbf{13.} $Re = 10^3\times 10^{-4}\times 10^{-5}/10^{-3} = 10^{-3}$.
\textbf{14.} $m = 1050\times 5.24\times 10^{-16} = \qty{5.5e-13}{kg}$;
$6\pi\eta r = 6\pi\times 10^{-3}\times 5\times 10^{-6} =
\qty{9.4e-8}{kg/s}$; $d = 5.5\times 10^{-13}\times 10^{-4}/9.4\times
10^{-8} = \qty{5.8e-10}{m}$: setengah nanometer.
\textbf{15.} $v_s = 2\times 25\times 10^{-12}\times 50\times
9.81/(9\times 10^{-3}) = \qty{2.7e-6}{m/s}$, \qty{2.7}{\micro m/s}.
\textbf{16.} Tenggelam \qty{1}{m}: $1/2.7\times 10^{-6} =
\qty{3.7e5}{s}$, 4,3 hari; berenang naik pada \qty{100}{\micro m/s}:
\qty{1e4}{s}, 2,8 jam. Berenang menang dengan faktor 37.
\textbf{17.} $c_\infty = \qty{1e-2}{mol/m^3}$: $J = 4\pi\times
1.9\times 10^{-9}\times 5\times 10^{-6}\times 10^{-2} =
\qty{1.2e-15}{mol/s} = 7.2\times 10^8$ molekul tiap sekon.
\textbf{18.} \qty{100}{pg} $= \qty{1e-10}{g} = \qty{8.3e-12}{mol}$
karbon, berlipat dua dalam \qty{12}{h} $= \qty{43200}{s}$:
\qty{1.9e-16}{mol/s} $= 1.2\times 10^8$ atom tiap sekon. Pasokan
dibagi permintaan: $7.2/1.2 = 6$.
\textbf{19.} Pasokannya tumbuh sepuluh kali ($\propto r$),
permintaannya seribu kali ($\propto r^3$): rasionya menjadi
$6/100 = 0.06$; sel semacam itu hanya dapat tumbuh pada enam persen
lajunya, sehingga alga yang disuapi difusi seukuran itu mustahil
tanpa pengadukan, pengangkutan internal, atau metabolisme yang jauh
lebih lambat.
\textbf{20.} $\Delta\Pi = RT\Delta c = 8.314\times 293\times 99 =
\qty{2.4e5}{Pa}$, yaitu \qty{2.4}{bar}.
\textbf{21.} Luasnya $4\pi(5\times 10^{-6})^2 = \qty{3.14e-10}{m^2}$;
aliran masuknya $L_p A\Delta\Pi = 10^{-14}\times 3.14\times 10^{-10}\times
2.4\times 10^5 = \qty{7.6e-19}{m^3/s} = \qty{0.76}{\micro m^3/s}$.
\textbf{22.} $524/0.76 = \qty{690}{s}$, kira-kira sebelas menit.
\textbf{23.} $P = \Delta\Pi\times$ aliran $= 2.4\times 10^5\times
7.6\times 10^{-19} = \qty{1.8e-13}{W}$.
\textbf{24.} Satu ATP: $5\times 10^4/6.02\times 10^{23} =
\qty{8.3e-20}{J}$: $1.8\times 10^{-13}/8.3\times 10^{-20} = 2.2\times
10^6$ ATP tiap sekon. Fiksasi karbonnya: $3\times 1.2\times 10^8 =
3.6\times 10^8$ ATP tiap sekon. Vakuolanya memakan \qty{0.6}{\%} dari
jumlah itu --- asuransi yang murah terhadap pecah.
\textbf{25.} Aliran air masuk \qty{0.76}{\micro m^3/s} (volume selnya
dalam sebelas menit); tanpa \omterm{prop:b2:unicellular-diversity:osmoregulation}{vakuola kontraktil} selnya akan
melipatduakan volumenya dalam kira-kira \qty{700}{s}; tetap kering
memakan sedikitnya $2\times 10^6$ ATP tiap sekon, kurang dari satu
persen dari yang dibelanjakan fotosintesisnya untuk karbon.
\end{solution}
